%!TEX root = ../thesis.tex
\chapter{Conclusions and Future Work}
\label{chap:conclusions}

% ============================================================================
\section{Summary of Contributions}
\label{sec:concl-summary}

This thesis evolved through three linked questions rather than a single benchmark objective. It began from the hypothesis that self-supervised backbones should improve even lightweight end-to-end planners. The publication of Drive-JEPA~\cite{wang-drive-jepa} confirmed that JEPA-style pretraining could deliver very strong mixed-city performance, which pushed the thesis toward a sharper question: what does SSL buy beyond pooled leaderboard PDMS? The answer developed here is cross-city robustness, together with the representational mechanism behind it.

The contributions below are those of the thesis. LAW~\cite{li-law-2024} is used as an external open-loop baseline, not claimed as an original method; the pretrained SSL checkpoints used throughout are adopted from prior work. The evaluation framework, experiments, analyses, and conclusions built around those components are original.

\begin{enumerate}
    \item \textbf{Cross-city evaluation framework and benchmark.} A per-city partitioning of NAVSIM and a paired open-loop (nuScenes) and closed-loop (NAVSIM) city-disjoint protocol. The partition first appears in our companion cross-city study~\cite{naeinian-cross-city-2026}; this thesis extends the protocol with a NAVSIM v1.1 re-evaluation and a broader representation-centred investigation (Chapter~\ref{chap:experiments}).

    \item \textbf{Cross-city SSL generalisation}. Under the evaluated settings with LAW open-loop on nuScenes and TransFuser / Latent TransFuser closed-loop on NAVSIM, frozen self-supervised representations (I-JEPA, DINOv2, MAE) reduce cross-city generalisation gaps relative to supervised baselines. The load-bearing conclusion: \textbf{cross-city failure is strongly mediated by the visual representation.}

    \item \textbf{DiffusionDrive $+$ SSL}. SSL backbones are integrated into DiffusionDrive through the same feature adapter used in the cross-city benchmark, including a camera-only variant. The completed mixed-city results show that the representation signal survives a planner swap under pooled evaluation. At thesis freeze, three rows of the cross-city DiffusionDrive matrix are also complete, including a matched Pittsburgh pair in which frozen I-JEPA shows a $2.0$-point smaller Singapore drop than ResNet-34. The full four-city matrix remains incomplete, so the generative cross-city result is treated as supportive evidence rather than as the primary thesis claim.

    \item \textbf{Mechanistic analysis}. A linear-probe, CKA, MMD, and feature-covariance effective-rank study of the trained checkpoint roster. The three calibrated findings are: (a) the linear-probe ordering on matched Boston-only Latent TransFuser rows (\reftab{tab:concl-probe}) tracks the out-of-distribution PDMS drop ordering; (b) fine-tuning compresses feature-covariance rank for every self-supervised family; (c) the signal is strongest within matched comparisons and attenuates when backbone scale, pretraining data, and training city vary together. No universal-mechanism claim is made.

\begin{table}[ht]
\centering
\caption[City-ID probe accuracy vs.\ OOD PDMS drop.]{City-ID linear-probe accuracy (second column) and Boston-only Latent TransFuser OOD PDMS drop in percentage points (third column) for the four backbone families. The rows are ordered by probe accuracy; the OOD-drop column is the corresponding value read from~\reftab{tab:navsim_pdms_boston_v11}. The two columns track each other: less city-decodable features transfer with a smaller geographic penalty.}
\label{tab:concl-probe}
\small
\setlength{\tabcolsep}{6pt}
\renewcommand{\arraystretch}{1.15}
\begin{tabular}{@{}l c c@{}}
\toprule
\textbf{Backbone} & \textbf{Probe acc.} & \textbf{OOD drop (pp)} \\
\midrule
ResNet-34 (supervised) & $0.97$ & $16.1$ \\
MAE                    & $0.89$ & $13.2$ \\
I-JEPA                 & $0.83$ & $9.6$ \\
DINOv2                 & $0.81$ & $6.3$ \\
\bottomrule
\end{tabular}
\end{table}
\end{enumerate}

% ============================================================================
\section{Key Findings}
\label{sec:concl-findings}

The project-level takeaway is that lightweight-head experiments motivated the representation hypothesis, Drive-JEPA demonstrated the in-distribution ceiling on pooled NAVSIM, and zero-shot cross-city transfer turns out to be the regime where representation choice matters most.

Four operational conclusions follow directly from the evidence in Chapter~\ref{chap:experiments}.

\begin{enumerate}
    \item \textbf{Representation choice is the largest single lever for cross-city robustness.} In the evaluated settings, swapping a supervised backbone for a frozen SSL ViT reduces open-loop transfer inflation by nearly an order of magnitude and improves closed-loop out-of-distribution PDMS by several points, with the planner held fixed~\cite{naeinian-cross-city-2026}.

    \item \textbf{The in-distribution cost is small; the out-of-distribution gain is large.} The DiffusionDrive all-city table shows at most a small mixed-city cost for SSL, while the cross-city LAW and TransFuser / Latent TransFuser matrices show a materially larger transfer gain under explicit geographic shift. The completed cross-city DiffusionDrive rows point in the same direction on the matched Pittsburgh pair, but the full four-city matrix is not yet complete.

    \item \textbf{Frozen SSL backbones generalise better than fully fine-tuned ones in this setting.} Fine-tuning improves mixed-city PDMS on several backbones, but compresses feature-covariance rank and, on the matched Boston-only Latent TransFuser rows, drifts the representation toward city-specific cues. Freezing is the operational lever; it is not a claim that fine-tuning is always harmful.

    \item \textbf{In the evaluated settings, representation choice mediates transfer more strongly than head choice.} The same ordinal pattern appears under LAW open-loop on nuScenes and TransFuser / Latent TransFuser closed-loop on NAVSIM, and the DiffusionDrive integration preserves the representation-dependent structure under both pooled evaluation and the completed cross-city rows. This is the basis for the bounded version of the central claim. The stronger claim, that the pattern is fully planner-invariant across all source cities, still depends on a completed cross-city DiffusionDrive matrix.
\end{enumerate}

% ============================================================================
\section{Industrial Implications}
\label{sec:concl-industrial}

The practical implication of the thesis is straightforward. A strong pooled benchmark score does not guarantee that an end-to-end planner will remain reliable when the operating city changes. For deployment teams, this means that geographic holdout evaluation should be treated as part of model validation rather than as an optional stress test.

The completed results also suggest a concrete design preference. When transfer to a new city is likely, a frozen self-supervised backbone is often the safer operating point than a supervised baseline that is slightly stronger on the pooled benchmark. The transfer gains are larger than the in-distribution losses, and the frozen setting avoids the feature-rank collapse that appears under aggressive fine-tuning.

There is also a systems implication. The camera-only DiffusionDrive results show that stronger image representations help lower-cost sensing stacks, but they do not erase the value of explicit geometry. This matters for industrial planning because it separates two questions that are often conflated: which representation transfers better, and which sensor suite is sufficient for a target deployment budget. The thesis answers the first question directly and gives a bounded partial answer to the second.

% ============================================================================
\section{Limitations}
\label{sec:concl-limitations}

\begin{itemize}
    \item \textbf{Benchmark-metric evaluation.} The thesis uses benchmark metrics (LAW open-loop ratios and NAVSIM PDMS) rather than an interactive closed-loop simulator. A simulator such as CARLA or nuPlan would provide stronger evidence of deployment-time generalization.

    \item \textbf{Limited geographic diversity.} NAVSIM/nuPlan covers only 4 cities (3 US + Singapore). Broader evaluation across more diverse geographies (Europe, developing countries, varied climate zones) is needed to fully validate the SSL generalization claim.

    \item \textbf{ImageNet pre-training.} Most SSL backbones were pre-trained on ImageNet, not driving data. While our nuScenes-pretrained I-JEPA shows competitive results, a comprehensive study of pre-training data domain is left for future work.

    \item \textbf{No weather/lighting variation.} Our cross-city protocol isolates geographic shift but does not independently test weather or lighting robustness, which is a complementary domain shift axis.

    \item \textbf{The full generative cross-city matrix is incomplete.} Three cross-city DiffusionDrive rows are available and directionally consistent with the main benchmark, but the full four-city matched matrix is not complete at thesis freeze. The thesis therefore does not claim planner-invariance under city-disjoint generative transfer.

    \item \textbf{Camera-only constraint.} While motivated by edge deployment, removing LiDAR sacrifices 3D geometry information that may be critical in certain scenarios (e.g., occluded pedestrians, narrow alleys).
\end{itemize}

% ============================================================================
\section{Future Work}
\label{sec:concl-future}

\subsection{Video SSL Backbones (V-JEPA)}
The most immediate next step is to move from image-based SSL to video-based SSL. Driving is inherently temporal, so a predictive video objective should align more naturally with planning than single-frame pretraining. Drive-JEPA already suggests that V-JEPA-style pretraining can be especially strong in-distribution, which makes it a natural candidate for the next iteration of the cross-city robustness question.

\subsection{Closed-Loop Evaluation}
Closed-loop evaluation is the most important methodological extension. Although PDMS is stronger than raw displacement error, it remains a benchmark proxy metric. A closed-loop simulator such as CARLA or the nuPlan simulator would make it possible to test whether the same representation trends hold when the ego policy can alter future scene evolution.

\subsection{Driving-Specific Pre-Training at Scale}
Most backbones in the thesis still inherit generic large-scale pretraining. A larger study of driving-specific pretraining at scale would clarify how much of the cross-city robustness gain comes from the SSL objective itself and how much comes from domain alignment between pretraining and deployment.

\subsection{Scaling to More Cities and Datasets}
The current city-split benchmark covers only four cities and one data pipeline. Extending the same protocol to other datasets, and especially to cross-dataset transfer, would test whether the representation findings generalize beyond NAVSIM and nuScenes.

\subsection{Multi-Modal JEPA Encoders}
A shared latent-space pretraining objective over camera and LiDAR would preserve the SSL robustness benefits while recovering part of the geometric signal lost in camera-only deployment. This is the direct extension suggested by the camera-only DiffusionDrive results in Chapter~\ref{chap:experiments}.
